\section{Revisão Bibliográfica}
\label{sec:revisao}

\subsection{Ergonomia no ambiente de trabalho}

A ergonomia é uma área multidisciplinar que tem como objetivo adaptar as condições
de trabalho às características físicas, cognitivas e psicossociais dos indivíduos,
promovendo conforto, segurança e eficiência. Segundo Iida e Guimarães (2016), a
ergonomia busca otimizar a interação entre o ser humano e os sistemas de trabalho,
reduzindo riscos à saúde e aumentando a produtividade.

No contexto de atividades de escritório, a ergonomia assume um papel ainda mais
relevante, uma vez que grande parte das tarefas é realizada em posição sentada e com
uso contínuo de computadores. De acordo com Dul e Weerdmeester (2012), a postura
adotada durante o trabalho influencia diretamente o desempenho e o bem-estar do
indivíduo, sendo fundamental a adequação do mobiliário e dos equipamentos utilizados.

A postura correta em ambientes de escritório envolve o alinhamento adequado da
coluna vertebral, apoio dos pés no chão, posicionamento correto dos braços e ajuste
da altura do monitor à linha dos olhos. Grandjean (1998) destaca que a manutenção de
uma postura inadequada por longos períodos pode resultar em sobrecarga
musculoesquelética, favorecendo o surgimento de dores e lesões.

Além disso, fatores como iluminação, organização do espaço e pausas regulares também
fazem parte das boas práticas ergonômicas. A ausência desses cuidados pode
comprometer não apenas a saúde física, mas também a concentração e o rendimento do
trabalhador.

\subsection{Distúrbios osteomusculares relacionados ao trabalho (DORT)}

Os distúrbios osteomusculares relacionados ao trabalho (DORT), também conhecidos
como \emph{Work-Related Musculoskeletal Disorders} (WMSDs), representam um dos
principais problemas de saúde ocupacional na atualidade. Esses distúrbios afetam
estruturas como músculos, tendões, articulações e nervos, sendo causados
principalmente por esforços repetitivos, posturas inadequadas e sobrecarga física.

Estudos indicam que a prevalência desses distúrbios entre trabalhadores de
escritório é elevada. Okezue et al. (2020) identificaram que aproximadamente
71,9\,\% dos trabalhadores analisados apresentavam algum tipo de DORT, sendo as
regiões mais afetadas a lombar, ombros e membros superiores. De forma semelhante,
Mohammadian et al. (2025) apontam uma prevalência de 80,81\,\%, com maior incidência
de dores no pescoço, coluna lombar e ombros.

Esses dados evidenciam a magnitude do problema, que não apenas compromete a saúde
dos trabalhadores, mas também gera impactos econômicos significativos, como aumento
do absenteísmo e redução da produtividade. Além disso, os DORT podem evoluir para
condições crônicas, limitando a capacidade funcional do indivíduo e afetando sua
qualidade de vida.

Os principais fatores de risco associados a esses distúrbios incluem posturas
inadequadas, permanência prolongada em posição estática, mobiliário inadequado,
repetitividade de movimentos e ausência de pausas durante a jornada de trabalho.
Fatores psicossociais, como estresse e pressão por produtividade, também contribuem
para o agravamento dessas condições (MOHAMMADIAN et al., 2025).

\subsection{Métodos de avaliação ergonômica}

A avaliação ergonômica é fundamental para identificar riscos e propor melhorias nas
condições de trabalho. Diversos métodos foram desenvolvidos com o objetivo de
analisar posturas, movimentos e organização do ambiente laboral.

Um dos instrumentos mais utilizados é o \emph{Nordic Musculoskeletal Questionnaire},
desenvolvido por Kuorinka et al. (1987), que permite identificar a presença de
sintomas musculoesqueléticos em diferentes regiões do corpo, sendo amplamente
aplicado em estudos epidemiológicos.

Outro método relevante é o RULA (\emph{Rapid Upper Limb Assessment}), proposto por
McAtamney e Corlett (1993), que avalia o risco de lesões nos membros superiores a
partir da análise da postura adotada durante a execução de tarefas. Esse método
considera fatores como posição dos braços, pescoço, tronco e força aplicada.

Além disso, o método ROSA (\emph{Rapid Office Strain Assessment}), apresentado por
Sonne, Villalta e Andrews (2012), é voltado especificamente para ambientes de escritório,
permitindo avaliar a ergonomia de estações de trabalho com base em critérios como
altura da cadeira, posicionamento do monitor e uso de periféricos.

Essas ferramentas são essenciais para a identificação de inadequações ergonômicas e
para a implementação de medidas corretivas, servindo também como base para o
desenvolvimento de sistemas automatizados de monitoramento postural. Em particular,
RULA e ROSA pontuam o risco a partir de \textbf{ângulos de inclinação de tronco e
pescoço}, o que fundamenta a escolha de pitch/roll do IMU como variável de decisão
deste trabalho.

\subsection{Sistemas inteligentes aplicados à ergonomia}

Com o avanço da tecnologia, os sistemas inteligentes têm sido cada vez mais
utilizados na área da ergonomia, permitindo o monitoramento contínuo e em tempo real
das condições de trabalho. Esses sistemas utilizam sensores, dispositivos
vestíveis, visão computacional e algoritmos de análise de dados para identificar
padrões de comportamento e detectar posturas inadequadas.

Um exemplo relevante é o sistema BackWatch, desenvolvido por Ferreira, Souza e Cubas
(2025), que utiliza sensores integrados para monitorar a inclinação do tronco e da
cabeça do usuário. O sistema é capaz de identificar desvios posturais e fornecer
alertas ao usuário, promovendo a autocorreção e prevenindo problemas de saúde. No
protótipo, dois sensores ultrassônicos de distância instalados no encosto da
cadeira, na altura da cabeça/pescoço e do quadril/cóccix, medem o afastamento do
usuário em relação ao encosto; um afastamento superior a cerca de 11\,cm em
qualquer um deles é tratado como postura inadequada e sinalizado em um display LCD
(alerta visual). Os próprios autores apontam, como trabalhos futuros, a troca do
display por um LED e um aplicativo de smartphone.

Além disso, estudos como o de Robertson et al. (2013) demonstram que intervenções
baseadas em ergonomia, aliadas ao uso de tecnologias, podem reduzir
significativamente os sintomas musculoesqueléticos e melhorar o desempenho dos
trabalhadores.

Os sistemas inteligentes apresentam vantagens como monitoramento contínuo,
personalização das recomendações e maior engajamento do usuário. No entanto, ainda
existem desafios relacionados à precisão dos sensores, adaptação a diferentes
biotipos e custo de implementação. \textbf{A proposta deste trabalho diferencia-se do
BackWatch ao fundir a pressão na cadeira (FSR) com a orientação do tronco superior
(IMU vestível)} --- combinação que o BackWatch, baseado apenas na distância ao
encosto com um limiar fixo, não realiza. As subseções seguintes revisam os sensores
empregados (Seção~\ref{sec:rev-sensores}), os trabalhos relacionados
(Seção~\ref{sec:rev-relacionados}), o uso de feedback háptico para a correção
postural (Seção~\ref{sec:rev-feedback}) e os limiares angulares usados na
literatura (Seção~\ref{sec:rev-limiares}), que fundamentam as escolhas do sistema
integrado descrito na Seção~\ref{sec:metodologia}.

% ---------------------------------------------------------------------------
\subsection{Sensores para monitoramento postural}
\label{sec:rev-sensores}

O monitoramento da postura sentada pode ser feito por sensores no mobiliário, que
medem a interação do corpo com a superfície de contato, ou por sensores vestíveis,
que medem a orientação de segmentos do corpo. Esta seção descreve as duas famílias
usadas neste trabalho e as limitações que condicionam o seu uso.

\subsubsection{Resistores sensíveis à força (FSR)}
\label{sec:rev-fsr}

Os resistores sensíveis à força (\emph{Force Sensing Resistors}, FSR) são
dispositivos de filme espesso polimérico (\emph{polymer thick film}, PTF) cuja
resistência elétrica diminui à medida que aumenta a força aplicada à sua superfície
(INTERLINK ELECTRONICS, [2012?a]). O efeito é piezorresistivo: o compósito de
polímero com partículas condutoras tem a sua condução entre partículas alterada
pela compressão, e o comportamento reológico (viscoelástico) do polímero introduz
uma dependência temporal da resistência sob carga constante (PAREDES-MADRID et al.,
2017). O modelo FSR402 usado neste trabalho pertence à série FSR 400 da Interlink
Electronics, com área ativa circular de 12{,}7\,mm de diâmetro (INTERLINK
ELECTRONICS, [2012?b]).

Segundo o fabricante, a série FSR 400 tem força de acionamento de cerca de
0{,}2\,N, faixa de sensibilidade de aproximadamente 0{,}2 a 20\,N, resistência sem
carga superior a 10\,M$\Omega$, repetibilidade de $\pm$2\,\% em uma mesma peça e de
$\pm$6\,\% entre peças de um mesmo lote, histerese média de +10\,\% e deriva de
longo prazo inferior a 5\,\% por década de tempo (ensaio de 35 dias sob 1\,kg)
(INTERLINK ELECTRONICS, [2012?a]). A curva força--resistência é fortemente não
linear e, por conveniência, é apresentada pelo fabricante em escalas logarítmicas
nos dois eixos (INTERLINK ELECTRONICS, [2012?b]). O guia
de integração do fabricante é explícito ao afirmar que o FSR \emph{não} é uma
célula de carga, um extensômetro ou um transdutor de pressão: em medições dinâmicas
obtêm-se, em geral, apenas resultados qualitativos, e a exatidão em força varia de
cerca de $\pm$6\,\% a $\pm$50\,\% conforme a consistência do sistema mecânico de
acionamento e do ambiente (INTERLINK ELECTRONICS, [2012?b]). O mesmo guia recomenda
o circuito divisor de tensão com um resistor de medição escolhido para maximizar a
sensibilidade na faixa de interesse e limitar a corrente, e alerta que a leitura
estática depende de um tempo de acomodação mecânica da ordem de segundos.

A deriva sob carga constante (\emph{creep}) é a limitação mais relevante para uma
pessoa sentada por longos períodos. Paredes-Madrid et al. (2017), estudando 32
amostras de FSR, mostraram que a resposta de \emph{creep} afeta fortemente a
exatidão das leituras de força e depende da tensão de alimentação do sensor. Por
essas razões, a literatura e o fabricante convergem para o uso do FSR como medida
\textbf{relativa} de distribuição de carga --- comparar sensores entre si, ou um
sensor com a sua própria leitura de referência --- e não como medida absoluta de
força. É essa a forma de uso adotada neste trabalho: a regra de decisão trabalha
com somas de grupos de sensores normalizadas por uma linha de base capturada na
calibração de cada usuário (Seção~\ref{sec:met-regra}).

\subsubsection{Matrizes e redes de sensores de pressão em assentos}
\label{sec:rev-matrizes}

Os primeiros trabalhos de cadeiras instrumentadas usaram mapas de pressão de alta
resolução, tratados como imagens em tons de cinza (TAN; SLIVOVSKY; PENTLAND, 2001).
Mapas desse tipo exigem tapetes comerciais de custo elevado, o que motivou duas
linhas de simplificação: reduzir o número de sensores discretos, escolhendo as
posições mais informativas (MUTLU et al., 2007), ou substituir o tapete por
sensores mais baratos e confortáveis, como sensores capacitivos têxteis (MEYER et
al., 2010) e matrizes têxteis resistivas com calibração própria (XU et al., 2013).
Meyer et al. (2010) observam que o material compressível do sensor têxtil apresenta
histerese significativa; modelando-a, o erro médio de medição caiu de 24\,\% para
5\,\%. Abordagens ainda mais enxutas usam células de carga sob o assento (ROH et
al., 2018) ou FSRs discretos distribuídos em assento, encosto e braços (ZEMP et al.,
2016).

Em revisão sistemática sobre redes de sensores para monitoramento da postura
sentada, Kappattanavar et al. (2021) identificaram que 13 dos 14 estudos incluídos
usavam sensores de pressão, e apenas 2 usavam acelerômetro --- um deles fixado no
encosto da cadeira, e não no corpo. Os autores recomendam incorporar IMUs para
capturar o movimento da coluna e as transições entre posturas, e destacam a falta
de definições padronizadas de postura como obstáculo à comparação entre estudos.

\subsubsection{Acelerômetros e IMUs para inclinação do tronco}
\label{sec:rev-imu}

Uma unidade de medição inercial (IMU) combina acelerômetro e giroscópio (e, em
alguns modelos, magnetômetro). O acelerômetro mede a diferença entre a aceleração
linear do sensor e o campo gravitacional local; na ausência de aceleração linear,
a sua saída é o vetor gravidade rotacionado para o referencial do sensor, do qual
se obtêm diretamente os ângulos de arfagem (\emph{pitch}) e rolagem (\emph{roll})
(PEDLEY, 2013). Essa estimativa estática tem três limitações conhecidas: (i) depende
da hipótese de aceleração linear desprezível, de modo que qualquer aceleração
dinâmica (movimento, impacto, vibração) contamina o ângulo; (ii) o ângulo de
guinada (\emph{yaw}), rotação em torno da vertical, não é observável pela
gravidade; e (iii) a solução se torna instável perto de $\pm$90\degree{} de
arfagem (PEDLEY, 2013).

O giroscópio, por sua vez, mede velocidade angular e não sofre com a aceleração
linear, mas a integração da sua saída acumula erro (deriva) ao longo do tempo.
Por isso, para movimentos dinâmicos, a orientação de segmentos corporais costuma
ser estimada por fusão dos dois sensores --- por exemplo, com filtro de Kalman que
usa o acelerômetro para corrigir continuamente a deriva da integração do
giroscópio (LUINGE; VELTINK, 2005). Wong e Wong (2008) aplicaram módulos com
acelerômetro triaxial e giroscópio no tronco superior, no tronco médio e na pelve
para estimar a inclinação do tronco e a variação da curvatura da coluna,
validando-os contra um sistema óptico de captura de movimento (Vicon 370) com nove
sujeitos.

A postura sentada é, em grande parte do tempo, quase estática, situação em que a
estimativa por gravidade é adequada e dispensa o giroscópio. Este trabalho usa
apenas o acelerômetro do LSM6DS3TR-C (giroscópio desligado), com
\emph{pitch} e \emph{roll} calculados pelas equações~(\ref{eq:met-pitch})
e~(\ref{eq:met-roll}). A limitação (i) aparece concretamente no sistema: a vibração
do próprio motor de alerta contamina a leitura, o que é tratado no firmware
(Seção~\ref{sec:met-corpo}). Além disso, a leitura absoluta depende de como a placa
assenta nas costas a cada uso, o que motiva a calibração por usuário em vez de um
limiar absoluto (Seção~\ref{sec:met-regra}).

% ---------------------------------------------------------------------------
\subsection{Trabalhos relacionados}
\label{sec:rev-relacionados}

\paragraph{Cadeiras instrumentadas.} Tan, Slivovsky e Pentland (2001) instalaram
sensores de distribuição de pressão no assento e no encosto de uma cadeira de
escritório e, aplicando análise de componentes principais aos mapas de pressão,
classificaram posturas estáticas com 96\,\% de acerto para usuários conhecidos e
79\,\% para usuários novos. Mutlu et al. (2007) propuseram uma estratégia de
posicionamento quase ótimo de sensores para reduzir custo e complexidade: obtiveram
87\,\% de acerto para dez posturas com 31 sensores em validação cruzada e 78\,\%
com 19 sensores em uma implantação física. Zemp et al. (2016) instrumentaram
cadeiras de escritório com 16 sinais de FSR (Interlink FSR 406: dez no assento,
quatro no encosto e três em cada braço) e um módulo inercial no encosto, usado para
estimar o ângulo do encosto; com 41 sujeitos e sete posições, o classificador
\emph{Random Forest} atingiu 90{,}9\,\% de acerto médio em sujeitos não vistos no
treinamento. Roh et al. (2018) mostraram que quatro células de carga sob o assento
bastam para classificar seis posturas (97{,}2\,\% de acerto médio, SVM, nove
sujeitos). Em todos esses trabalhos a saída é a \emph{classificação} da postura,
sem retorno ao usuário.

Outros trabalhos acrescentam feedback à cadeira. Haller et al. (2011) usaram uma
cadeira com quatro sensores para comparar alertas gráficos, físicos (movimento da
cadeira) e vibrotáteis. Zheng e Morrell (2013) desenvolveram uma cadeira de
orientação postural que compara a postura atual com uma postura de referência e
avalia feedback vibrotátil e visual. O LifeChair (BOURAHMOUNE; AMAGASA, 2019) é uma
almofada com sensores de pressão, aplicativo de smartphone e quatro atuadores
vibratórios, com acerto de 98{,}93\,\% em mais de 13 posturas; os autores destacam a
influência do índice de massa corporal do usuário no padrão de pressão.

\paragraph{Vestíveis com sensores inerciais.} Wong e Wong (2008), já citados,
observaram de forma preliminar que os sujeitos melhoravam a postura quando
recebiam sinais de feedback do sistema vestível. Breen, Nisar e ÓLaighin (2009)
usaram um único acelerômetro para estimar o ângulo craniovertebral e fornecer
biofeedback visual em tempo real a usuários de computador. Kuo et al. (2021)
avaliaram um sensor vestível que vibra quando o usuário curva o tronco
(Seção~\ref{sec:rev-feedback}). O BackWatch (FERREIRA; SOUZA; CUBAS, 2025), embora
monitore a inclinação do tronco e da cabeça, é um sistema montado na cadeira
(sensores ultrassônicos), e não vestível.

\paragraph{Fusão de modalidades.} A combinação de pressão no assento com orientação
do tronco medida no corpo é pouco explorada. Entre os 14 estudos revisados por
Kappattanavar et al. (2021), nenhum combinava sensores de pressão na cadeira com
IMU vestível; o caso mais próximo é o de Zemp et al. (2016), cujo módulo inercial
está no encosto e mede a inclinação da cadeira, não a do usuário. A
Tabela~\ref{tab:rev-comparativo} resume os trabalhos discutidos.

\begin{table}[H]
\centering
\caption{Comparação entre trabalhos relacionados e o sistema proposto.}
\label{tab:rev-comparativo}
\begin{singlespace}
\footnotesize
\begin{tabularx}{\textwidth}{>{\raggedright\arraybackslash}p{2.7cm}
                             >{\raggedright\arraybackslash}X
                             >{\raggedright\arraybackslash}p{2.1cm}
                             >{\raggedright\arraybackslash}X
                             >{\raggedright\arraybackslash}p{2.0cm}}
\toprule
\textbf{Trabalho} & \textbf{Sensores} & \textbf{Local} &
\textbf{Saída / feedback} & \textbf{Calibração por usuário} \\
\midrule
Tan, Slivovsky e Pentland (2001) & Mapas de pressão de alta resolução &
  Assento e encosto & Classificação (96\,\% / 79\,\%); sem feedback & --- \\
Mutlu et al. (2007) & 19--31 sensores de pressão (posicionamento otimizado) &
  Cadeira & Classificação de 10 posturas (87\,\% / 78\,\%); sem feedback & --- \\
Zemp et al. (2016) & 16 sinais de FSR 406 + módulo inercial &
  Assento, encosto, braços; IMU no encosto & Classificação de 7 posições
  (90{,}9\,\%); sem feedback & Não (modelo geral) \\
Roh et al. (2018) & 4 células de carga & Sob o assento &
  Classificação de 6 posturas (97{,}2\,\%) & --- \\
Haller et al. (2011) & 4 sensores na cadeira & Cadeira &
  Alerta gráfico, físico ou vibrotátil & --- \\
Zheng e Morrell (2013) & Sensores na cadeira & Cadeira &
  Feedback vibrotátil ou visual & Postura de referência \\
Bourahmoune e Amagasa (2019) & Sensores de pressão (almofada) &
  Assento & Classificação ($>$13 posturas) + vibração + aplicativo & --- \\
Wong e Wong (2008) & 3 módulos (acelerômetro + giroscópio) &
  Corpo: tronco e pelve & Inclinação do tronco + sinais de feedback & --- \\
Breen, Nisar e ÓLaighin (2009) & 1 acelerômetro & Corpo: cabeça/pescoço &
  Biofeedback visual em tempo real & --- \\
Kuo et al. (2021) & Sensor vestível de biofeedback & Corpo (vestível) &
  Vibração ao curvar o tronco & --- \\
Ferreira, Souza e Cubas (2025) & 2 sensores ultrassônicos de distância &
  Encosto & Alerta visual (LCD), limiar fixo de distância & Não \\
\midrule
\textbf{Este trabalho} & 8 FSR402 + acelerômetro (IMU LSM6DS3TR-C) &
  Assento e encosto + corpo (alto das costas) & Alerta sim/não decidido na
  cadeira; vibração no vestível & Sim (botão: referência do IMU + linha de base
  dos FSRs) \\
\bottomrule
\end{tabularx}
\end{singlespace}
\fonte{Elaborado pelo autor (2026), com base nas fontes citadas. Percentuais =
acerto de classificação relatado pelos autores (usuários conhecidos / novos, ou
validação cruzada / implantação), obtidos com protocolos diferentes e, portanto,
não comparáveis entre si. ``---'' = não descrito nas fontes consultadas.}
\end{table}

\paragraph{Posicionamento deste trabalho.} Diante desse panorama, o sistema
proposto diferencia-se em quatro aspectos. Primeiro, \textbf{funde as duas
modalidades}: a cadeira observa como o peso se distribui (assento traseiro
descarregado, deslocamento lateral, cadeira vazia) e o vestível observa a
orientação do tronco superior, a parte do corpo que não toca a cadeira; cada canal
cobre casos que o outro confunde, como reclinar no encosto (ângulo grande, mas
postura aceitável) e escorregar no assento (tronco pouco inclinado, mas apoio
traseiro perdido). Segundo, fornece \textbf{feedback háptico em tempo real} no
próprio corpo, em vez de apenas registrar ou classificar a postura. Terceiro, usa
\textbf{calibração por usuário}, capturando por botão a referência ``ereto'' do IMU
e a linha de base dos FSRs na mesma janela, em vez de limiares absolutos ou de
modelos treinados em outras pessoas. Quarto, todo o \textbf{processamento é
embarcado} --- a cadeira decide o alerta e o devolve ao vestível pelo enlace sem fio
---, sem aplicativo de smartphone nem nuvem. O objetivo é \emph{detectar} postura
prejudicial (alerta sim/não), não maximizar o acerto de classificação; por isso,
e porque a validação deste trabalho usa amostra pequena, os percentuais da
Tabela~\ref{tab:rev-comparativo} não devem ser comparados diretamente com os
resultados da Seção~\ref{sec:resultados}.

% ---------------------------------------------------------------------------
\subsection{Feedback háptico e biofeedback para correção postural}
\label{sec:rev-feedback}

O biofeedback postural consiste em devolver ao usuário, em tempo real, uma
informação sobre a própria postura que ele não percebe espontaneamente durante a
tarefa, de modo que a correção parta dele. Há evidências de curto prazo de que
esse retorno induz autocorreção. Breen, Nisar e ÓLaighin (2009), com seis sujeitos
monitorados por cinco horas com e sem biofeedback visual, observaram redução
significativa da fração de tempo em postura inadequada em todos eles. Kuo et al.
(2021), em estudo quase-experimental com 21 adultos jovens que digitaram por uma
hora com e sem um sensor vestível que vibra quando o usuário curva o tronco,
mediram com captura de movimento tridimensional redução da flexão do pescoço
($p < 0{,}001$), da cifose torácica ($p = 0{,}033$) e do ângulo pélvico
($p = 0{,}021$) com o dispositivo; os autores ressaltam, porém, que se tratou de uma
única sessão.

Quanto à modalidade, Zheng e Morrell (2013) não encontraram diferença
estatisticamente significativa entre feedback vibrotátil e visual na capacidade de
levar o usuário a uma postura de referência; em tarefa dupla (digitar enquanto
recebia o feedback), ambos reduziram o desempenho de digitação em grau semelhante.
O argumento a favor do canal tátil é que, no trabalho de escritório, a visão já
está ocupada pela tela. A vibração, entretanto, tem custo: Haller et al. (2011)
verificaram que o feedback vibrotátil parecia o mais útil para chamar a atenção no
início, mas era o mais intrusivo, sendo o feedback físico (movimento da cadeira)
avaliado como o menos disruptivo ao trabalho.

A evidência de efeito sustentado é mais fraca. Em ensaio de viabilidade com 62
trabalhadores da área da saúde, Ribeiro et al. (2014) observaram redução da postura
lombar flexionada no grupo com feedback constante em relação à linha de base, mas
sem diferença estatisticamente significativa entre os grupos. No ensaio
randomizado por \emph{clusters} subsequente, controlado por dispositivo simulado
(\emph{sham}), o feedback sonoro não reduziu o número de vezes em que o limiar
postural foi excedido (RIBEIRO et al., 2020). Esses resultados indicam que o alerta
precisa ser pouco intrusivo, para não ser abandonado, e que o efeito em longo prazo
não pode ser presumido a partir de testes curtos.

No sistema proposto, essas considerações se traduzem em escolhas de projeto: o
alerta só dispara após o desvio persistir por 3\,s (evitando reagir a movimentos
breves), só termina abaixo de um limiar menor (histerese de 10\degree{} para
7\degree{}), não dispara com a cadeira vazia nem em posturas aceitáveis como
reclinar no encosto, e a vibração é aplicada no corpo, onde é percebida sem
desviar a visão da tela (Seção~\ref{sec:met-regra}). A adesão e o efeito em uso
prolongado não foram avaliados neste trabalho.

% ---------------------------------------------------------------------------
\subsection{Ângulos e limiares de postura sentada}
\label{sec:rev-limiares}

Os métodos observacionais pontuam o risco por faixas de ângulo absoluto em relação
à vertical. No RULA (MCATAMNEY; CORLETT, 1993), o pescoço recebe +1 para flexão de
0 a 10\degree{}, +2 de 10 a 20\degree{}, +3 acima de 20\degree{} e +4 em extensão;
o tronco recebe +1 quando ereto (inclusive sentado corretamente), +2 de 0 a 20\degree{} de
flexão, +3 de 20 a 60\degree{} e +4 acima de 60\degree{}, com +1 adicional em ambos
os casos para torção ou inclinação lateral. No ROSA (SONNE; VILLALTA; ANDREWS, 2012),
o encosto recebe a menor pontuação quando oferece apoio lombar adequado com
reclinação entre 95 e 110\degree{}, e é penalizado quando o ângulo fica abaixo de
95\degree{} ou acima de 110\degree{}. A norma ISO 11226 (INTERNATIONAL ORGANIZATION
FOR STANDARDIZATION, 2000) estabelece limites recomendados para posturas de
trabalho estáticas combinando ângulos corporais e tempo de manutenção da postura.

Dois pontos dessas referências orientam o sistema. Primeiro, o risco depende tanto
do ângulo quanto do \textbf{tempo} de permanência, o que justifica exigir que o
desvio persista antes de alertar. Segundo, as faixas do RULA para o pescoço têm
largura de 10\degree{}, mesma ordem do limiar de 10\degree{} usado neste trabalho.
A comparação, porém, é apenas indicativa: o sistema mede o \textbf{desvio} em
relação à referência ``ereto'' de cada usuário, com o sensor no alto das costas, e
não o ângulo absoluto de um segmento em relação à vertical; além disso, o
reclinar no encosto, que o ROSA considera aceitável dentro da faixa recomendada, é
tratado pela fusão com os FSRs, e não pelo ângulo. Os valores de 10\degree{},
7\degree{} e 3\,s vieram dos dados deste trabalho
(Seção~\ref{sec:met-regra}), não de uma norma.
